\section[Cote 161-2 : monadicité, descente, caractérisation de Gabriel--Ulmer]{Cote 161-2 : monadicité, descente, caractérisation de Gabriel--Ulmer\protect\verifie{Folder161_2}}\label{sec:161-2}

La cote 161-2, \emph{Algèbre universelle [ou catégories] : notes manuscrites (s.d.)}, cent onze pages que l'inventaire date \q{[vers 1963-1973]}, cherche ce qu'est une espèce de structure algébrique dans une catégorie et comment les modèles déterminent la théorie. Deux feuilles de travail, pages~33 et~34, notent le critère de reconnaissance d'un foncteur \q{tripleable}, avec trois noms que la transcription lit avec doute, \q{Beck, Barr, Linton} (\lot{161-2}{2}{33}), puis le critère de descente pour un morphisme de topos, conclu par \q{$B_\pi \xrightarrow{\ \simeq\ } B$ (th)} (\lot{161-2}{2}{34}). La page~60 énonce un corollaire sur les catégories cocomplètes, \q{Ce sont ces $E$ qu'il faudrait appeler « accessibles »} (\lot{161-2}{3}{60}), que la page~61 démontre (\lot{161-2}{4}{61}). La page~77 établit, pour les théories définies par limites finies, le lemme $\mathrm{Ind}(\Sigma) \simeq S$ (\lot{161-2}{4}{77}).

\noindent\emph{Conventions.} Les notations de la section~\ref{sec:19} valent ici : pour une adjonction $f \dashv g$, $f : A \to B$, $h : A \to \varphi\text{-}\mathbf{Coalg}$ est le foncteur de comparaison vers les coalgèbres de $\varphi = fg$ ; dualement, pour $F \dashv U$, $U : A \to B$, le foncteur de comparaison va de $A$ dans les algèbres de la monade $\mathbb T = UF$. Une double flèche est \emph{réflexive} si ses deux flèches ont une section commune, \emph{coréflexive} si elles ont une rétraction commune.

\begin{enonce}[Théorème de monadicité, page 33]
Un foncteur $U : A \to B$ admettant un adjoint à gauche $F$ est monadique ($A \simeq \mathrm{Alg}(\mathbb T)$, $\mathbb T = UF$) si $U$ est conservatif et si $A$ a et $U$ préserve les coégalisateurs des paires réflexives.
\end{enonce}

\begin{enonce}[Critère de descente, page 34]
Soit $f$ un morphisme de topos, d'image inverse $f^* : B' \to B$, et $\pi = f^* f_*$. Si $f^*$ est conservatif, le foncteur de comparaison $B' \to B_\pi$, $X \mapsto f^* X$, vers la catégorie des $\pi$-coalgèbres, est une équivalence.
\end{enonce}

\begin{lemme}\label{lem:161-2-grossier}
Soit $f \dashv g$ une adjonction, $f : A \to B$. Si $f$ est conservatif, si $A$ a les noyaux des doubles flèches coréflexives et si $f$ y commute, le foncteur de comparaison $h$ est une équivalence.
\end{lemme}

\begin{proof}
Pour toute coalgèbre $(X, a)$, la double flèche $(\eta_{gX}, g(a))$ est coréflexive (lemme~\ref{lem:19-paire}) ; elle a donc un noyau, auquel $f$ commute. Le lemme~\ref{lem:19-comp}\,(3) conclut. C'est le théorème de comonadicité « grossier » de Beck, classique.
\end{proof}

\begin{proof}[Démonstration du théorème de monadicité]
C'est l'énoncé dual. Le foncteur $U^\circ : A^\circ \to B^\circ$ est adjoint à gauche de $F^\circ$ ; la monade $UF$ sur $B$ devient la comonade $U^\circ F^\circ$ sur $B^\circ$, dont les coalgèbres sont les opposées des $\mathbb T$-algèbres, et le foncteur de comparaison devient celui du lemme~\ref{lem:161-2-grossier}. Les coégalisateurs de paires réflexives de $A$ sont les noyaux de paires coréflexives de $A^\circ$, et $U^\circ$ est conservatif. Le lemme~\ref{lem:161-2-grossier} s'applique.
\end{proof}

\begin{proof}[Démonstration du critère de descente]
On démontre davantage : soit $f^* \dashv f_*$ une adjonction quelconque, $f^* : B' \to B$, où $B'$ a les limites finies et $f^*$ est exact à gauche et conservatif. Alors $B'$ a les noyaux de toutes les doubles flèches, en particulier des coréflexives, et $f^*$, exact à gauche, y commute. Le lemme~\ref{lem:161-2-grossier} donne l'équivalence $B' \simeq B_\pi$. Pour un morphisme de topos, $B'$ a les limites finies et $f^*$ est exact à gauche par définition.
\end{proof}

\begin{definition}
Soient $E$ une catégorie localement petite et $\kappa$ un cardinal régulier. Un objet $X$ de $E$ est \emph{$\kappa$-présentable} si $\operatorname{Hom}_E(X, -)$ commute aux limites inductives $\kappa$-filtrantes, et \emph{présentable} s'il l'est pour un $\kappa$ ; c'est l'\q{accessible} de la page. Une sous-catégorie pleine $\mathcal C$ de $E$ est \emph{dense} si, pour tout $X$, la flèche canonique $\varinjlim_{\mathcal C_{/X}} Y \to X$ est un isomorphisme ; c'est la condition~a) par laquelle la page~56 définit une sous-catégorie \q{strictement génératrice}. La catégorie $E$ est \emph{localement présentable} si elle est cocomplète et s'il existe un cardinal régulier $\kappa$ et un petit ensemble d'objets $\kappa$-présentables dont tout objet est limite inductive $\kappa$-filtrante ; c'est ce que dit \q{$E$ est équivalente à un $\mathrm{Ind}_\pi(\mathcal C)$, $\mathcal C$ petite}.
\end{definition}

\begin{enonce}[Corollaire, pages 60 et 61]
Une catégorie localement petite et cocomplète $E$ est localement présentable si et seulement si elle admet une petite sous-catégorie pleine dense formée d'objets présentables.
\end{enonce}

\begin{proof}
Supposons $E$ localement $\kappa$-présentable. La sous-catégorie pleine $P$ des objets $\kappa$-présentables est essentiellement petite et dense : c'est classique (Gabriel--Ulmer, Adámek--Rosický). Soit $Q \subseteq P$ une petite sous-catégorie pleine contenant un représentant de chaque classe d'isomorphisme ; l'inclusion $Q \to P$ est une équivalence, et le composé d'une équivalence et d'un foncteur dense est dense, donc $Q$ est dense dans $E$. Ses objets sont $\kappa$-présentables, donc présentables.

Réciproquement, soit $P$ une petite sous-catégorie pleine dense dont chaque objet $X$ est $\kappa_X$-présentable pour un cardinal régulier $\kappa_X$. Comme $P$ est petite, la famille des $\kappa_X$ est majorée par un cardinal $c$. Soit $\kappa_0$ le successeur de $\max(c, \aleph_0)$ ; c'est un cardinal régulier, comme tout successeur d'un cardinal infini. Toute limite inductive $\kappa_0$-filtrante étant $\kappa_X$-filtrante, chaque $X \in P$ est $\kappa_0$-présentable. Une sous-catégorie dense est fortement génératrice (classique). Or une catégorie cocomplète qui admet un petit ensemble fortement générateur d'objets $\kappa_0$-présentables est localement $\kappa_0$-présentable : c'est le théorème~1.20 d'Adámek et Rosický, classique. Donc $E$ est localement présentable.
\end{proof}

\begin{enonce}[Lemme, page 77]
Soient $R$ une petite catégorie à limites finies, $S = \underline{\mathrm{Hom}}_{\mathrm{lex}}(R, \mathrm{Ens})$ et $\Sigma \simeq R^\circ \subset S$ la sous-catégorie pleine des structures représentables. Le foncteur d'inclusion $\Sigma \to S$ induit une équivalence $\mathrm{Ind}(\Sigma) \simeq S$.
\end{enonce}

\begin{proof}
On identifie $S$ aux préfaisceaux $X : \Sigma^\circ \to \mathrm{Ens}$ exacts à gauche, la catégorie $\Sigma = R^\circ$ ayant les limites inductives finies, et $\mathrm{Ind}(\Sigma)$ à la sous-catégorie pleine des préfaisceaux sur $\Sigma$ qui sont limites inductives filtrantes de représentables. Il suffit de montrer que ces deux sous-catégories pleines ont les mêmes objets. Un représentable est exact à gauche, et les limites inductives filtrantes commutent aux limites projectives finies dans $\mathrm{Ens}$ ; tout ind-objet est donc exact à gauche. Réciproquement, soit $X$ exact à gauche. Tout préfaisceau est la limite inductive des représentables sur sa catégorie des éléments $\Sigma_{/X}$, et celle-ci est filtrante : elle est non vide, car $X$ envoie l'objet initial de $\Sigma$ sur un point ; deux éléments $x \in X(a)$, $y \in X(b)$ proviennent de $X(a \sqcup b) = X(a) \times X(b)$ ; deux flèches $u, v : (a, x) \rightrightarrows (b, y)$ sont égalisées par la flèche vers le coégalisateur $c$ de $u, v$, puisque $X(c)$ est le noyau de $X(u), X(v) : X(b) \rightrightarrows X(a)$, qui contient $y$. C'est le fait classique que les ind-objets d'une petite catégorie à limites inductives finies sont ses préfaisceaux exacts à gauche (SGA~4, exposé~I, §~8), que la preuve vérifiée invoque.
\end{proof}

\begin{remarque}
Le critère de descente n'utilise pas la structure de topos. Il vaut pour tout foncteur adjoint à gauche, exact à gauche et conservatif, défini sur une catégorie à limites finies. L'équivalence a pour source $B'$, alors que la page écrit $B$ ; la lecture l'a corrigé.
\end{remarque}

\begin{remarque}
Le corollaire est le théorème de caractérisation des catégories localement présentables de Gabriel et Ulmer. Il est démontré ici pour la notion moderne ; le $\mathrm{Ind}_\pi(\mathcal C)$ de SGA~4 que la page emploie n'est pas formalisé. Le lemme de la page~77 est le cas des limites finies de la dualité de Gabriel--Ulmer.
\end{remarque}
